\documentclass[10pt,a4paper]{amsart}
\usepackage[margin=19mm]{geometry}
\usepackage{amsmath,amssymb,mathtools}
\usepackage[scr=rsfs,cal=euler]{mathalfa}
\usepackage{xfrac}
\usepackage[unicode,hidelinks,bookmarks=false]{hyperref}
\newcommand{\ie}{i.e.\ }
\newcommand{\cf}{cf.\ }
\newcommand{\resp}{respectively\ }
\newcommand{\Subsection}{\subsection}
\providecommand{\nobreakdash}{\nobreak\hbox{-}}
\setlength{\emergencystretch}{2em}
\multlinegap0pt
\usepackage{amssymb}
\usepackage{float}
\usepackage{xcolor}
\newtheorem{proposition}{Proposition}
\providecommand{\llbracket}{\mathopen{[\mkern-2.7mu[}}
\providecommand{\rrbracket}{\mathclose{]\mkern-2.7mu]}}
\title{Spectral partitioning for $k$-block averaging kernels of finite Markov chains}
\author{Michael C.H. Choi, Youjia Wang}
\date{Source: arXiv:2608.21466v1 (CC BY 4.0)}
\begin{document}
\maketitle

\noindent\emph{Source and licence.} Spectral partitioning for $k$-block averaging kernels of finite Markov chains. Version arXiv:2608.21466v1 (CC BY 4.0), distributed by arXiv under the Creative Commons Attribution 4.0 International licence, http://creativecommons.org/licenses/by/4.0/, as recorded in the arXiv licence metadata for this exact version and shown on the abstract page https://arxiv.org/abs/2608.21466v1. Full licence notice: \texttt{licenses/arxiv-2608.21466-CC-BY-4.0.txt}.\par\smallskip

\noindent\textit{Editorial scope.} Counted unit: the article's proposition prop:additive-mixture-identity (source0.tex lines 1342--1377). The article's complete original proof of it is delivered with it (source0.tex lines 1379--1455, one \texttt{proof} environment of 77 lines). No mathematics is written, repaired or re-derived: the statement and the proof are the article's own lines, in the article's own notation, and no retained line is substituted or re-flowed.  No further source-local context is needed: every cross-reference of every retained line resolves inside the two delivered spans.  Nothing was omitted from any retained span, so the package carries no omission record.  No retained line contains a citation, so the wrapper carries no bibliography and no bibliography entry.  No retained line contains a figure environment, an image-inclusion command or drawing code, so no image is delivered and the package declares no asset dependency.  Editorial formatting only, in addition: an AMS \texttt{amsart} preamble that restates the article's own declarations and macros used by the retained lines, namely the environments \texttt{proposition} on separate counters, together with the standard packages those lines need.  The retained environments print in this document as Proposition 1 in order of appearance, whereas the article prints the same items with the numbering of its own full document.  The complete native source of the article as distributed in the arXiv e-print for this exact version is bundled unchanged as \texttt{sources/208251/source0.tex}.
\begin{proposition}[Additive-mixture identity]
	\label{prop:additive-mixture-identity}
	Let $k\in\llbracket 2,n-1\rrbracket$.
	For every $k$-block partition $\mathcal O$,
	\begin{equation}
		H(\mathcal O)
		=\operatorname{Tr}\left(
		\frac{P^2}{4}
		+\frac{G_{\mathcal O}P}{2}
		+\frac{G_{\mathcal O}}{4}
		-\Pi
		\right).
		\label{eq:H-trace}
	\end{equation}
	Moreover,
	\begin{equation}
		H(\mathcal O)
		=C_{P,k}
		+\frac12\operatorname{Tr}(Q_{\mathcal O}P),
		\label{eq:H-reduction}
	\end{equation}
	where
	\begin{equation}
		C_{P,k}
		:=\frac14\operatorname{Tr}(P^2)+\frac{k-2}{4}.
		\label{eq:H-constant}
	\end{equation}
	Consequently, for fixed $P$ and $k$,
	\begin{equation*}
		\operatorname*{argmin}_{|\mathcal O|=k}H(\mathcal O)
		=
		\operatorname*{argmin}_{|\mathcal O|=k}
		\operatorname{Tr}(Q_{\mathcal O}P).
		%\label{eq:H-equivalent-objective}
	\end{equation*}
\end{proposition}

\begin{proof}
	Since $P$ and $G_{\mathcal O}$ are self-adjoint,
	$A_{\mathcal O}$ is self-adjoint.  Moreover,
	\[
	A_{\mathcal O}\Pi
	=\Pi A_{\mathcal O}
	=\Pi.
	\]
	It follows that
	\[
	\begin{aligned}
		H(\mathcal O)
		&=\operatorname{Tr}\left(
		(A_{\mathcal O}-\Pi)^*
		(A_{\mathcal O}-\Pi)
		\right) \\
		&=\operatorname{Tr}\left(
		(A_{\mathcal O}-\Pi)^2
		\right) \\
		&=\operatorname{Tr}(A_{\mathcal O}^2-\Pi).
	\end{aligned}
	\]
	Because
	\[
	A_{\mathcal O}
	=\frac{P+G_{\mathcal O}}{2}
	\]
	and $G_{\mathcal O}^2=G_{\mathcal O}$, we have
	\[
	A_{\mathcal O}^2
	=\frac14\left(
	P^2+PG_{\mathcal O}
	+G_{\mathcal O}P+G_{\mathcal O}
	\right).
	\]
	Cyclicity of the trace gives
	\[
	\operatorname{Tr}(PG_{\mathcal O})
	=\operatorname{Tr}(G_{\mathcal O}P).
	\]
	Substituting into the preceding expression proves
	\eqref{eq:H-trace}.  Notice that this is an equality of traces;
	the corresponding operators need not be equal because $P$ and
	$G_{\mathcal O}$ need not commute.
	
	It remains to isolate the partition-dependent term.  Since
	\[
	G_{\mathcal O}=\Pi+Q_{\mathcal O},
	\]
	we have
	\[
	\operatorname{Tr}(G_{\mathcal O}P)
	=\operatorname{Tr}(\Pi P)
	+\operatorname{Tr}(Q_{\mathcal O}P)
	=1+\operatorname{Tr}(Q_{\mathcal O}P).
	\]
	Furthermore, $G_{\mathcal O}$ is a rank-$k$ orthogonal projection, so
	\[
	\operatorname{Tr}(G_{\mathcal O})=k.
	\]
	Therefore, \eqref{eq:H-trace} becomes
	\[
	\begin{aligned}
		H(\mathcal O)
		&=\frac14\operatorname{Tr}(P^2)
		+\frac12\left(
		1+\operatorname{Tr}(Q_{\mathcal O}P)
		\right)
		+\frac{k}{4}-1 \\
		&=\frac14\operatorname{Tr}(P^2)
		+\frac{k-2}{4}
		+\frac12\operatorname{Tr}(Q_{\mathcal O}P),
	\end{aligned}
	\]
	which proves \eqref{eq:H-reduction}.  The first two terms depend only
	on $P$ and $k$, so the equality of the two sets of minimizers follows.
\end{proof}

\end{document}
